% ==============================================================================
% CAPÍTULO 10: MARKETING (Orquestador del Capítulo)
% ==============================================================================
\preportadacapitulo{10}{Marketing}

\newpage
\section{Capítulo 10: Marketing}\label{cap:marketing}

% [INSTRUCCIÓN]: Marketing digital. Contenido propio de esta versión.
% [PENDIENTE]: montos del presupuesto (10.1.4.4 y 10.1.5) y las cuentas reales en redes: no se inventan; se completan con cotizaciones de las plataformas y con la decisión del equipo.

Este capítulo plantea cómo NEXWARE da a conocer el sistema entre las instituciones que podrían adoptarlo. Como el cliente es una institución (Capítulo~\ref{cap:mercado}), el marketing es de tipo consultivo: pocos contactos, de alto valor, en los que importan la confianza y la demostración del sistema más que el volumen de visitas.

\subsection{10.1. Marketing Digital}

\subsubsection{10.1.1. Análisis de la situación}

\paragraph{10.1.1.1. Fortalezas del producto para la comunicación.}

Los puntos del producto que mejor se comunican son concretos y medibles:

\begin{itemize}
    \item Se evalúa el 100\,\% de las solicitudes, no solo el 40\,\% que hoy se revisa (sección 1.2).
    \item Los cupos se asignan por necesidad y con trazabilidad, no por azar ni por orden de llegada.
    \item El estudiante postula y consulta su estado desde el celular.
    \item El sistema reduce el papel y el trabajo del personal.
    \item Hay un caso real de uso: la DUBSS de la UAGRM.
\end{itemize}

\paragraph{10.1.1.2. Público objetivo segmentado.}

Se distinguen dos públicos con mensajes distintos:

\begin{itemize}
    \item \textbf{Quien decide la compra:} autoridades universitarias, directivos de bienestar estudiantil y responsables de sistemas. Les interesa la transparencia, el ahorro de tiempo y el cumplimiento de las auditorías.
    \item \textbf{Quien usa el sistema:} estudiantes postulantes y personal administrativo. Les interesa que el trámite sea sencillo y justo.
\end{itemize}

\paragraph{10.1.1.3. Posicionamiento del producto.}

NEXWARE se posiciona como la empresa local que entiende el proceso de la beca comedor y lo convierte en un sistema propio de cada institución: «Si no existe, lo programamos.» (sección 2.7). El mensaje central es que el sistema se hizo con la DUBSS, para su proceso real y no a partir de una plantilla.

\subsubsection{10.1.2. Estrategia de contenidos especializados}

\paragraph{10.1.2.1. Contenido educativo.}

Guías breves sobre cómo digitalizar el proceso de una beca: de qué datos depende el puntaje, cómo se prepara el expediente digital y cómo se audita un finalista. Se publican en el sitio de NEXWARE (Capítulo~\ref{cap:sitio_web}).

\paragraph{10.1.2.2. Contenido informativo.}

Avisos para los estudiantes: fechas de la convocatoria, documentos necesarios, cómo consultar el estado de la postulación y cómo instalar la aplicación. Este contenido lo publica la institución cliente en sus propios canales y NEXWARE aporta el material gráfico.

\paragraph{10.1.2.3. Contenido institucional.}

Presentación de NEXWARE: quiénes la forman, su misión, visión y eslogan, y la forma de trabajo con pruebas automatizadas y control de versiones (Capítulo~\ref{cap:sqap}).

\paragraph{10.1.2.4. Casos de uso y demostraciones.}

Videos cortos y capturas de pantalla que muestran el recorrido de un postulante y de un auditor con datos de prueba, nunca con datos reales de estudiantes. El resumen del caso de la DUBSS se publica cuando la institución lo autorice por escrito.

\subsubsection{10.1.3. Canales de adquisición}

\paragraph{10.1.3.1. Facebook.}

Se usa para difundir contenido informativo y las novedades del sistema.

\paragraph{10.1.3.2. Instagram.}

Se usa para contenido visual breve: capturas de pantalla, publicaciones de «cómo se usa» y avisos de las fechas de la convocatoria.

\paragraph{10.1.3.3. LinkedIn.}

Es el canal para el público que decide la compra: se publica el contenido educativo e institucional y se contacta a responsables de otras instituciones.

\paragraph{10.1.3.4. TikTok.}

Se usa para videos cortos dirigidos a los estudiantes: cómo postular en pocos pasos y cómo revisar el estado de la solicitud.

\paragraph{10.1.3.5. Página web corporativa.}

Es el canal principal. Reúne la información de la empresa, la descripción del sistema y el formulario de contacto (sección 7.1). Todos los demás canales llevan a ella.

\subsubsection{10.1.4. Plan de publicidad digital}

La publicidad pagada se usa de forma limitada y solo después de tener un caso en operación que mostrar. Antes de ese momento se emplea contenido orgánico y contacto directo.

\paragraph{10.1.4.1. Publicidad en Meta Ads.}

Anuncios en Facebook e Instagram dirigidos a estudiantes de la institución cliente durante la convocatoria, para recordar fechas y requisitos.

\paragraph{10.1.4.2. Publicidad en LinkedIn Ads.}

Anuncios segmentados por cargo (directivos, responsables de bienestar estudiantil y de sistemas de instituciones de educación superior) para presentar el sistema.

\paragraph{10.1.4.3. Publicidad en TikTok Ads.}

Videos promocionados para estudiantes durante la convocatoria. Es el canal de menor prioridad para NEXWARE como vendedor.

\paragraph{10.1.4.4. Objetivos publicitarios.}

Los objetivos buscan resultados medibles: conseguir reuniones de demostración con instituciones y aumentar la cantidad de estudiantes que postulan por el sistema.

% [PENDIENTE]: metas numéricas (reuniones, alcance, costo por contacto). Definirlas con el equipo; no se inventan.

\subsubsection{10.1.5. Presupuesto de marketing}

% [PENDIENTE]: el presupuesto necesita las tarifas reales de cada plataforma y la decisión del equipo sobre cuánto invertir. Los cinco puntos siguientes quedan sin montos.

\paragraph{10.1.5.1. Gestión de redes sociales.}

El equipo gestiona las cuentas con sus propios integrantes; el costo es el tiempo de trabajo. % [PENDIENTE]: horas por semana y costo.

\paragraph{10.1.5.2. Publicidad digital.}

% [PENDIENTE]: monto mensual por plataforma (Meta, LinkedIn, TikTok).

\paragraph{10.1.5.3. Inversión mensual total.}

% [PENDIENTE]: suma de 10.1.5.1 y 10.1.5.2.

\paragraph{10.1.5.4. Inversión anual.}

% [PENDIENTE]: inversión mensual por los meses de campaña.

\paragraph{10.1.5.5. Justificación de la inversión.}

La inversión se justifica si el ingreso de un solo contrato (licencia y servicios, sección 8.7) supera el gasto de marketing del año. El monto que cubre ese criterio se calcula cuando se definan los precios.

\subsubsection{10.1.6. Anexos}

Las publicaciones, los afiches y los videos que se produzcan se guardan como anexos de este capítulo. Hasta ahora no se ha producido material.

% [PENDIENTE]: incorporar piezas gráficas reales (afiche de convocatoria, publicación de ejemplo) cuando existan.
